树图之外，“OPE 成立”在含夸克的 QCD 中是三件不同的事：欧氏非例外动量下 Green 函数的渐近展开；物理矩阵元中只出现规范不变算符；前向在壳部分子矩阵元（例外动量）也可以这样展开。下面把它们分开陈述，各自给出证明的位置与地位。

\begin{proposition}[欧氏非例外动量下 Green 函数的逐图渐近展开 \tagA]\label{prop:dis-ope-euclid}
考虑含 $n_f$ 个无质量夸克的 QCD，协变规范，规范参数 $\xi$ 属于含 $1$ 的紧集 $\Xi$，维数正规化加 $\overline{\rm MS}$，算符基取 $D$ 维基（含规范变算符、鬼场算符与 evanescent 算符）。设 $\Phi$ 是基本场的串，$(p,q)$ 为非例外的欧氏动量（外动量的任何真部分和都不为零），$\tilde G^T_{jj;\Phi}(q;p)$ 是两个电磁流的编时乘积与 $\Phi$ 的截腿 Green 函数，$\tilde G_{[O];\Phi}(p)$ 是重整化算符 $[O]$ 的插入。在微扰论的每一阶 $r$：
\begin{enumerate}[label=(A\arabic*),leftmargin=2.6em,itemsep=0pt,topsep=2pt]
\item $\tilde G^T_{jj;\Phi}-\sum_{d_O\le a}C_O(q;\mu,\xi)\,\tilde G_{[O];\Phi}=\order{|q|^{1-a}\ln^k|q|}$，对 $q$ 的方向、对 $p$ 属于非例外紧集、对 $\xi\in\Xi$ 一致；
\item $C_O=|q|^{2-d_O}\times\big(L\equiv\ln(q^2/\mu^2)$ 的至多 $r$ 次多项式$\big)$，系数是 $\xi$ 的多项式的协变张量；
\item $C^{\overline{\rm MS}}=C^{\rm bare}Z^{-1}$ 有限，其中 $C^{\rm bare}_O$ 是含两个流的硬子图的维数正规化 Taylor 系数之和，只依赖于 $(q,g_0,\xi_0,\varepsilon)$；
\item 在满足条件 (NE)（四维树图顶点线性无关）的每个块中系数唯一（\cref{lem:appM-uniqueness}(b)、\cref{lem:appM-ne}）；特别地，全部扭度 2 系数对一切 $n$ 唯一。
\end{enumerate}
（本命题只涉及部分子态的 Green 函数。）
\uses{lem:dis-tree-ope,thm:qcd-renormalizability,thm:ee-euclid-ir,lem:appM-uniqueness,lem:appM-ne}
% errtype: E-P
\end{proposition}

\begin{proofsketch}
(1) (A1)、(A3)：对每个图、在每个 $\xi$ 下引用大动量极限的渐近展开定理\cite{Smirnov:1990rz,Smirnov2002}与 $\overline{\rm MS}$ 中的 OPE\cite{Chetyrkin:1982zq}，再用 $O_{\rm bare}=Z^{-1}[O]$；非例外动量下维数正规化 Green 函数的红外有限性见\cref{thm:ee-euclid-ir}（\cite[定理~1]{Breitenlohner:1976te}）。$\xi$ 的一致性来自固定阶时余项是 $\xi$ 的多项式（Lagrange 插值）。$\Phi\ne\varnothing$ 时 $\tilde G^T$ 不需要相加的重整化：唯一的零量纲定域项 $g^{\mu\nu}AA$、$g^{\mu\nu}\bar cc$ 违反中性外腿的电磁 Ward 恒等式。(2) (A2)：在固定 $\varepsilon$ 下，$k$ 圈硬部分对 $(q;g_0,\xi_0)$ 是 $2-d_O-2k\varepsilon$ 次齐次的；代入 $g_0=\mu^\varepsilon Z_gg$、$\xi_0=Z_3\xi$ 与 $Z^{-1}$ 后为 $|q|^{2-d_O}\sum_kc_k(\varepsilon,\xi)(\mu^2/q^2)^{k\varepsilon}$，极点至多 $r$ 阶，有限的极限必是 $L$ 的多项式。(3) (A4)：\cref{lem:appM-uniqueness}(b) 与\cref{lem:appM-ne}。树图层次即\cref{lem:dis-tree-ope}。完整论证（含所引定理的适用范围）见附录~M 的\cref{sec:appM-massless}。（此步给出的是 E-P 型陈述，在所引定理适用之处逐图成立。）
\end{proofsketch}

\paragraph{地位与缺口}
标签为 [A]：所引的定理覆盖不了 (A1)、(A3) 的完整陈述。Smirnov 的定理（\cite[第~117~页]{Smirnov:1990rz}）是一个混合陈述——在坐标变量上是分布、在非例外欧氏动量上是函数，并且用的是解析正规化；维数正规化下的函数形式只是被断言（同文第~118~页）；鬼场、规范变算符基与复合旁观者都没有处理。另一条路线——用已发表的动量空间 OPE 定理加上系数唯一性去推出\cref{asm:appM-remainder}——是循环的（\cref{rem:appM-remainder-reduction}），而维数正规化路线根本不需要这一假设。误差类型为 E-P（逐图、逐点于 $q$）。升级的途径：一个覆盖含夸克、鬼场与规范变算符的 QCD、以维数正规化函数形式陈述的渐近展开定理（第~25~章列为开放问题）。

\begin{proposition}[物理矩阵元只需规范不变算符 \tagA]\label{prop:dis-ope-gi}
(a) 在物理部分子态（在壳，且动量转移非零或外腿红外正规化）之间，$j^\mu j^\nu$ 的 OPE 只有规范不变算符贡献；在扭度 2 只有 $O^{(n)}_q$、$O^{(n)}_g$。
\par
(b)（蕴含关系，\tagR，$R_{\mathrm{cond}}$）（$R_{\mathrm{cond}}$：在所列编号假设之下证明完整；结论的地位不高于其中最弱的假设。）前提：\cref{prop:dis-ope-euclid}、\cref{asm:appM-brst-ope}。若带一个复合旁观者 $Z=s\phi$ 的欧氏展开（含硬子图包含 $Z$ 顶点的合并项）成立，即\cref{asm:appM-brst-ope}的 (E$\varepsilon$)，则 (a) 成立：对每个量纲 $d$，系数和 $X_d=\sum_{d_O=d}C_O\,O$ 满足 $sX_d\in\mathcal E$（运动方程算符生成的理想），从而 $X_d$ 是规范不变算符、BRST 恰当算符与运动方程算符之和，后两类在物理态之间的矩阵元为零。
\par
（本命题只涉及部分子态。）
\uses{prop:dis-ope-euclid,thm:prelim-composite-renorm,lem:dis-twist-classification,asm:appM-brst-ope,thm:appM-local-cohomology,lem:appM-alien-vanish,lem:appM-uniqueness,thm:appM-brst-reduction,lem:appM-lbc-formal}
% errtype: E-P
\end{proposition}

\begin{proofsketch}
(b) 的完整证明见附录~M 的\cref{sec:appM-brst}（\cref{thm:appM-brst-reduction}），步骤为：(1) 维数正规化中 $sj^\mu=0$ 的 Ward 恒等式 $\sum_k\epsilon'_k\langle\mathcal J\phi_1\cdots(s\phi_k)\cdots\phi_n\rangle=0$；(2) 代入 (E$\varepsilon$)，非合并部分给出 $-\sum_OC_O\langle(sO)\Phi\rangle$；(3) 由维数正规化的 Schwinger--Dyson 恒等式（$\delta^{(D)}(0)=0$），合并部分组合为 $\langle\mathbb E(q)\Phi\rangle$，$\mathbb E\in\mathcal E$；(4) 二者之和对一切 $\Phi$ 都是 $\order{|q|^{1-a}\ln^k|q|}$，按量纲比较（\cref{lem:appM-uniqueness}(c)）迫使 $\sum_{d_O=d}C_O\,sO\in\mathcal E$；(5) 逐圈的定域 BRST 上同调（\cref{lem:appM-lbc-formal}，基于\cref{thm:appM-local-cohomology}）给出 $X_d=$ 规范不变 $+$ BRST 恰当 $+$ 运动方程算符；(6) 规范不变算符只混入规范不变算符与外来算符（\cref{thm:prelim-composite-renorm}），外来算符在物理态之间的矩阵元为零（\cref{lem:appM-alien-vanish}）；(7) 扭度 2 的规范不变算符由\cref{lem:dis-twist-classification}列出。(a) 由 (b) 与 (E$\varepsilon$) 得到。文献：Joglekar 处理了非阿贝尔理论中两个规范不变流的 OPE，并表明可能出现的规范变算符的物理矩阵元为零或不可分离（据该文摘要）\cite{Joglekar:1976eb,Joglekar:1976pe}；纯 Yang--Mills 中 OPE 在规范不变算符上闭合\cite{Hollands:2007zg}。
\end{proofsketch}

\paragraph{地位与缺口}
[A] 即\cref{asm:appM-brst-ope}中 (E$\varepsilon$) 的地位：带一个复合旁观者、含合并项的逐图维数正规化欧氏展开没有被证明。上述文献的结论与 (a) 一致，但它们对\cref{prop:dis-ope-euclid}中维数正规化 $\overline{\rm MS}$ 系数的适用性没有被确立（Joglekar 的处理在 QCD 流上，Hollands 的在纯 Yang--Mills 的 Epstein--Glaser 框架中）。前向在壳的注意事项：对零动量转移的在壳无质量胶子，外来算符矩阵元为零的论证不适用——Collins--Scalise 对一个 BRST 恰当算符给出了反例\cite{Collins:1994ee}，它解释了 Hamberg--van Neerven 发现的表观非三角混合\cite{Hamberg:1991qt}；对运动方程算符，前向在壳无质量的情形只是不被标准证明覆盖。误差类型 E-P。升级的途径：证明 (E$\varepsilon$)。

\begin{proposition}[前向在壳部分子矩阵元 \tagA]\label{prop:dis-ope-forward}
(a) 对纯维数正规化中在壳无质量的部分子靶 $j$（物理极化平均），在微扰论的每一阶，
\begin{equation}\label{eq:dis-forward-ope}
T_2^{(j)}(\omega)=4\sum_{n=2,4,\dots}\omega^{n-1}\sum_iC^{(n)}_{2,i}A^{(n)}_{i/j},
\end{equation}
在 $|\omega|\le r<1$ 上一致收敛；$A^{(n)}_{i/j}$ 是部分子矩阵元（$\delta_{ij}$ 加纯 $\overline{\rm MS}$ 抵消项）。
\par
(b)（蕴含关系，\tagR，$R_{\mathrm{cond}}$）（$R_{\mathrm{cond}}$：在所列编号假设之下证明完整；结论的地位不高于其中最弱的假设。）前提：\cref{prop:dis-ope-euclid}、\cref{prop:dis-ope-gi}、\cref{asm:appM-forward}。前向在壳动量是例外动量，\cref{prop:dis-ope-euclid}不能直接用；在前向在壳无质量部分子矩阵元的 OPE 余项估计（\cref{asm:appM-forward}）之下，(a) 成立。
\par
(c) 在每一阶，解析的一步是逐图证明了的：每个前向在壳维数正规化 Laurent 系数在割开的 $\omega$ 平面中全纯，并等于其朴素 Taylor 级数（\cref{thm:appM-forward-expansion}，\tagR，$R_{\mathrm{PT}}$）；(a) 作为\cref{cor:appM-forward-ope}(a) 成立（\tagA：与 $\overline{\rm MS}$ 系数的认同要经过\cref{prop:dis-ope-euclid}的 (A3) 与\cref{asm:appM-brst-ope}的 (E$\varepsilon$)）。
\par
（本命题只涉及部分子态。）
\uses{prop:dis-ope-euclid,prop:dis-ope-gi,lem:prelim-scaleless,asm:appM-forward,lem:appM-uniqueness,cor:appM-forward-ope,thm:appM-forward-expansion}
% errtype: E-P
\end{proposition}

\begin{proofsketch}
(b)（附录~M 的\cref{sec:appM-forward}）：(1) 在壳无质量部分子的算符矩阵元中一切圈积分都无标度（\cref{lem:prelim-scaleless}），故 $\langle p|[O]|p\rangle$ 只剩树图与纯极点；它只能是 $c\,p^{\mu_1}\cdots p^{\mu_n}$（$n$ 为自旋，含 $g^{\mu_i\mu_j}$ 的项正比于 $p^2=0$），而其量纲 $d_O-2$ 必须等于 $n$，否则无标度的 $c$ 只能为零，故只有扭度 $d_O-n=2$ 贡献。(2) 沿 $2p\cdot\hat q$ 为纯虚数的欧氏方向 $q=t\hat q$，$\omega=c_0/t$，由\cref{asm:appM-forward}余项为 $\order{t^{2-N}\ln^kt}$。(3) 若 $T_2^{(j)}$ 与\cref{eq:dis-forward-ope}之差在 $\omega^{N-3}$ 以下有非零的 Taylor 系数，它将给出 $\gtrsim t^{1-n_*}$ 的项，与 (2) 矛盾；$N$ 任意，故 (a) 成立。(c)：\cref{thm:appM-forward-expansion}的证明用 Lee--Pomeransky 参数表示、Newton 多面体控制与维数正规化的可容许性；\cref{cor:appM-forward-ope}再用\cref{lem:appM-uniqueness}、(A3) 与 (E$\varepsilon$) 把 Taylor 系数认同为 $\overline{\rm MS}$ 系数。单圈、$N\le8$ 的精确检验见\cref{sec:dis-wilson}。完整论证见附录~M 的\cref{sec:appM-forward}。
\end{proofsketch}

\paragraph{地位与缺口}
[A] 即\cref{asm:appM-forward}的地位。\cref{cor:appM-forward-ope}把它化归为 (A3)、(E$\varepsilon$) 与 Joglekar--Lee 三角性——是化归，不是解除。\cref{prop:dis-ope-gi}地位一段所述的前向反例针对 BRST 恰当算符，对运动方程算符前向在壳情形只是未被覆盖。因此不能说“前向在壳部分子矩阵元直接给出 Wilson 系数，无需额外步骤”：例外动量与外来算符正是需要额外步骤的地方。误差类型 E-P。升级的途径：解除\cref{cor:appM-forward-ope}的三个输入（第~25~章列为开放问题）。

\paragraph{振幅形式。}\cref{eq:dis-forward-ope}可以写成卷积：
\begin{equation}\label{eq:dis-amplitude-form}
T_2^{(j)}(\omega)=\sum_q\int_{-1}^1\dd\xi\,\hat H_q(\xi\omega)\,f_{q/j}(\xi)+\int_0^1\dd\xi\,\hat H_g(\xi\omega)\,f_{g/j}(\xi),\qquad \hat H_i(w)=4\sum_{n=2,4,\dots}C^{(n)}_{2,i}\,w^{n-1},
\end{equation}
其中 $f_{i/j}$ 是矩为 $\int_{-1}^1\xi^{n-1}f_{q/j}=A^{(n)}_{q/j}$、$\int_0^1\xi^{n-1}f_{g/j}=A^{(n)}_{g/j}$ 的分布（其算符定义见第~8~章）。夸克项按 $f_q(-\xi)=-f_{\bar q}(\xi)$ 奇延拓到 $[-1,0]$，于是 $\int_{-1}^1\xi^{n-1}f_q=\int_0^1\xi^{n-1}[q+(-1)^n\bar q]$；胶子项只取 $[0,1]$（$f_g(-\xi)=-f_g(\xi)$ 与 $\hat H_g$ 的奇性使 $[-1,1]$ 上的积分加倍）。树图 $\hat H_q(w)=4e_q^2w/(1-w^2)$，与\cref{ex:dis-free-quark}一致。式中没有 $\dd\xi/\xi$：$1/\xi$ 已被 $\tilde P^\mu\tilde P^\nu/(P\cdot q)$ 吸收。这是光锥展开在动量空间中的熟知形式（参见\cite{ChristHasslacherMueller1972,BrandtPreparata1971}）。

\begin{theorem}[重整化微扰论中的欧氏 OPE \tagA]\label{thm:dis-ope}
\cref{prop:dis-ope-euclid,prop:dis-ope-gi,prop:dis-ope-forward}同时成立。等价地：对带红外正规化部分子外动量的 $T\,j^\mu j^\nu$ 的 Green 函数，当 $Q^2\to\infty$（欧氏 $q$；对前向部分子矩阵元，则对 $\omega$ 的每个固定的 Taylor 系数），$T^{\mu\nu}$ 由\cref{eq:dis-ope-tensor}加上余项 $R$ 给出，求和遍及偶数 $n$ 与扭度 2 规范不变算符 $i$；系数 $C^{(n)}_{a,i}(Q^2/\mu^2,\as(\mu))$ 与外态及红外正规化无关；$R$ 是扭度 $\ge4$ 的贡献，对每个固定的 $n$ 为 $\order{Q^{-2}}$ 乘以对数。按\cref{def:dis-twist}归一化矩阵元时，$T_1\ni2\omega^nC^{(n)}_{1,i}A^{(n)}_i$，$T_2\ni4\omega^{n-1}C^{(n)}_{2,i}A^{(n)}_i$。% VF-DIS-13, VF-CONV-15
（本定理只涉及部分子态。）
\uses{prop:dis-ope-euclid,prop:dis-ope-gi,prop:dis-ope-forward}
% errtype: E-P
\end{theorem}

\begin{proofsketch}
由三个子命题合成：(1)\cref{prop:dis-ope-euclid}给出非例外欧氏动量下的展开与（逐块）唯一的系数；(2)\cref{prop:dis-ope-gi}把物理矩阵元中的算符限制为规范不变算符，扭度 2 为 $O_q$、$O_g$；(3)\cref{prop:dis-ope-forward}把展开推广到前向部分子矩阵元的每个固定 Taylor 系数。张量结构与收缩规则是\cref{eq:dis-ope-tensor,eq:dis-contraction}。奇扭度（手征奇）系数在无质量微扰论中为零，这一点在本节稍后单独证明，所以余项从扭度 4 开始。完整论证见附录~M 的\cref{sec:appM-massless}、\cref{sec:appM-brst}与\cref{sec:appM-forward}。
\end{proofsketch}

\paragraph{地位与缺口}
对标量场论，OPE 在重整化微扰论的每一阶被证明\cite{Zimmermann1973,Zimmermann:1972te,Zimmermann1970}（在 $\varphi^4$ 中甚至在每一固定圈数下在有限距离收敛\cite{Hollands:2011gf,Holland:2014pna}）；对纯 Yang--Mills 中的规范不变算符，OPE 在规范不变算符上闭合\cite{Hollands:2007zg}，红外的控制见\cite{Frob:2015uqy}。对含夸克的 QCD，Collins 的教科书断言这一展开 “is proved in a gauge theory the same way as in any other theory”\cite[第~12.8~节]{Collins1984}（在规范理论中与在其他理论中以同样方式被证明），IJJSZ 则写道它 “is widely assumed to hold for any renormalizable quantum field theory including QCD”\cite[第~II~节]{IzubuchiEtAl2018}（被广泛假定对包括 QCD 在内的一切可重整理论成立）；规范理论的必要成分见\cite{Joglekar:1976eb,Joglekar:1976pe}。对无质量可重整理论，Fröb--Holland 证明了 OPE 的存在，并在上同调条件 $H^{1,4}_{E(4)}(s|d)=0$ 下证明了系数的 BRST Ward 恒等式（预印本）\cite{Frob:2016mzv}；有限距离上的收敛只在 $\varphi^4_4$ 的固定圈数下被证明；Anikin--Zav'yalov 对“微扰论的任意阶”推导了短距离与光锥展开\cite{Anikin:1978tj}，其所覆盖的理论类没有被确认，不能当作 QCD 的证明引用。[A] 同时指三个子陈述：第一个因为所引定理的适用范围不足，后两个因为\cref{asm:appM-brst-ope}与\cref{asm:appM-forward}；前向部分子子陈述只被化归，没有被解除（\cref{cor:appM-forward-ope}）。在微扰论之外，本定理不能离开\cref{asm:intro-np-ope}与\cref{asm:intro-af-coeff}而成立。误差类型 E-P。

本定理与第~6~章真空关联函数的欧氏 OPE（\cref{thm:ee-ope}）步骤相同，只是外态现在是夸克与胶子、算符取一切自旋；文献中关于 OPE 严格性的其他结果汇总在第~24~章（\cref{rem:front-ope-qcd-status}）。

\begin{proposition}[奇扭度系数的手征压低 \tagR]\label{prop:dis-chiral-suppression}
（$R_{\mathrm{PT}}$：在重整化微扰论中对红外正规化的部分子态严格。）在无质量夸克的微扰论中，两个矢量流的编时乘积在夸克态之间的一切手征奇结构，在每一阶逐图为零：每条穿过硬部分的开放夸克线携带奇数个 $\gamma$ 矩阵。因此一切手征奇（由\cref{lem:dis-twist-chirality}，即奇扭度的对称无迹）算符的 OPE 系数在无质量微扰论中为零；有夸克质量时它们是 $m_q$ 的奇函数。（本命题只涉及部分子态。）
\uses{lem:dis-twist-chirality,def:prelim-feynman-rules}
\end{proposition}

\begin{proof}
无质量夸克的传播子为 $\ii\sla k/k^2$，夸克与胶子、夸克与光子的顶点为 $\ii g\gamma^\mu t^a$、$\ii e_q\gamma^\mu$；胶子与鬼场顶点不含 $\gamma$ 矩阵，闭合夸克圈给出迹，只贡献数值因子。一条从外夸克场到外夸克场、经过 $n_V$ 个矢量顶点的开放夸克线有 $n_V-1$ 个内传播子，其 Dirac 结构是 $2n_V-1$ 个 $\gamma$ 矩阵（$\gamma^\mu$ 或 $\sla k$）之积。Clifford 代数由 $\{\gamma^\mu,\gamma^\nu\}=2g^{\mu\nu}$ 定义，按 $\gamma$ 的个数是 $\mathbb Z_2$ 分次的，在任意 $D$ 中都成立；奇数个 $\gamma$ 之积是奇秩反对称积 $\gamma^{[r]}$（$r$ 奇）的线性组合。手征奇结构对应偶秩的 $\gamma^{[r]}$（$1$、$\sigma^{\mu\nu}$、……），其系数恒为零。同样地，OPE 系数乘在 $\bar\psi\,\Gamma_{\rm line}\psi$ 上，$\Gamma_{\rm line}$ 只含奇秩结构，故手征奇算符的系数为零。整个论证只用 $\gamma$ 的个数的奇偶性，不涉及 $\gamma_5$，因而在维数正规化中成立。夸克质量使传播子变为 $\ii(\sla k+m_q)/(k^2-m_q^2)$，偶秩结构只能与 $m_q$ 的奇次幂同时出现。
\end{proof}

非微扰地看，奇扭度系数的压低是\cref{asm:intro-np-ope}的 (O2) 的一部分：手征对称性被自发破缺，把全部味的 $\psi\to\gamma_5\psi$ 是转角 $\pi/2$ 的 $U(1)_A$ 转动，它把 $\theta$ 移动 $n_f\pi$；瞬子效应带有 $\ee^{-2\pi/\as(|q|)}$ 的压低。

\begin{proposition}[由部分子结构函数提取 Wilson 系数 \tagR]\label{prop:dis-wilson-partons}
（$R_{\mathrm{OPE}}$：以 OPE 为前提而严格。QCD 中含夸克的微扰 OPE 是全阶论证（\cref{thm:dis-ope}），严格证明只对标量理论与纯 Yang--Mills 的规范不变算符给出；对强子态还须非微扰 OPE（\cref{asm:intro-np-ope}）。）设 $j$ 是纯维数正规化中在壳无质量的部分子（物理极化平均），并设在所考虑的阶前向振幅的 Taylor 系数等于部分子结构函数矩的四倍：对偶数 $n\ge2$，$t^{(j)}_n=4\int_0^1\dd z\,z^{n-2}\hat F^{(j)}_2(z)$，$\tau^{(j)}_n=4\int_0^1\dd z\,z^{n-1}\hat F^{(j)}_1(z)$（矩恒等式；对红外正规化的靶，这是\cref{prop:dis-dispersion}在增长条件下的结论）。则对每个偶数 $n\ge2$，
\begin{equation}\label{eq:dis-wilson-partons}
\int_0^1\dd z\,z^{n-2}\hat F^{(j)}_2(z)=\sum_iC^{(n)}_{2,i}\,A^{(n)}_{i/j},\qquad A^{(n)}_{i/j}=\delta_{ij}+(\text{纯 }\overline{\rm MS}\text{ 抵消项}),
\end{equation}
对 $\hat F_L$ 与 $C^{(n)}_{L,i}=C^{(n)}_{2,i}-C^{(n)}_{1,i}$ 有同样的关系。（本命题只涉及部分子态。）
\uses{thm:dis-ope,prop:dis-ope-forward,lem:prelim-scaleless,prop:dis-dispersion}
\end{proposition}

\begin{proof}
由\cref{prop:dis-ope-forward}(a)（OPE 前提），$t^{(j)}_n=4\sum_iC^{(n)}_{2,i}A^{(n)}_{i/j}$；与假设 $t^{(j)}_n=4\int z^{n-2}\hat F^{(j)}_2$ 比较即得\cref{eq:dis-wilson-partons}的第一式。矩阵元：$\langle j|[O_i]|j\rangle=\sum_kZ_{ik}\langle j|O_k^{\rm bare}|j\rangle$。在零动量转移下外态只有动量 $p$，$p^2=0$，裸矩阵元的一切圈积分没有任何标度，由\cref{lem:prelim-scaleless}为零；树图矩阵元为 $\langle j|O_k|j\rangle_{\rm tree}=2\delta_{kj}(p^{\mu_1}\cdots p^{\mu_n})$（\cref{def:dis-twist}(a)(b)）。故 $A^{(n)}_{i/j}=Z_{ij}(n)$，在 $\overline{\rm MS}$ 中是 $\delta_{ij}$ 加纯极点；这些极点原本是紫外的，在无标度积分中与红外极点抵消后，留下的是红外（共线）极点的记账。对 $T_1$，由\cref{thm:dis-ope}的收缩规则 $\tau^{(j)}_n=2\sum_iC^{(n)}_{1,i}A^{(n)}_{i/j}$ 与假设得 $2\int z^{n-1}\hat F_1=\sum_iC^{(n)}_{1,i}A^{(n)}_{i/j}$，再与 $\hat F_L/z=\hat F_2/z-2\hat F_1$ 合起来即 $C_L$ 的关系。
\end{proof}

因此 Wilson 系数可以从部分子结构函数（割图）读出：在维数正规化中，部分子的 PDF（矩阵元）是纯抵消项，结构函数矩的有限部分就是 $\overline{\rm MS}$ 系数；PDF 的对应陈述是第~8~章的\cref{prop:pdf-partonic-dr}。单圈的精确检验（$N\le8$ 时矩恒等式精确到 $\varepsilon$ 成立）见\cref{ex:dis-oneloop-ope}。

\begin{proposition}[Wilson 系数与约化矩阵元的重整化群 \tagR]\label{prop:dis-ope-rg}
（$R_{\mathrm{OPE}}$：以 OPE 为前提。）记 $\Gamma_{ij}(n,\as)$ 为扭度 2 算符的反常维度矩阵：
\begin{equation}\label{eq:dis-rg}
\frac{\dd A^{(n)}_i}{\dd\ln\mu^2}=\sum_j\Gamma_{ij}(n,\as)A^{(n)}_j,\qquad \frac{\dd C^{(n)}_{a,j}}{\dd\ln\mu^2}=-\sum_iC^{(n)}_{a,i}\Gamma_{ij}(n,\as),\qquad \Gamma_{ij}=\frac{\as}{2\pi}\gamma_{ij}(n)+\order{\as^2},
\end{equation}
其中 $A^{(n)}_q$ 按味计、包括夸克与反夸克，$\Gamma_{q_kg}=2\gamma_{qg}$、$\Gamma_{gq_k}=\gamma_{gq}$，最低阶的值为% VF-CONV-16, VF-PDF-8, VF-PDF-10
\begin{gather}
\gamma_{qq}(n)=C_F\Big[\frac32+\frac1{n(n+1)}-2S_1(n)\Big],\qquad \gamma_{qg}(n)=T_F\frac{n^2+n+2}{n(n+1)(n+2)},\qquad \gamma_{gq}(n)=C_F\frac{n^2+n+2}{(n-1)n(n+1)},\nonumber\\
\gamma_{gg}(n)=2C_A\Big[\frac1{n(n-1)}+\frac1{(n+1)(n+2)}-S_1(n)\Big]+\frac{\beta_0}2,\label{eq:dis-gamma0}
\end{gather}
即领头阶 DGLAP 核的 Mellin 矩 $\gamma_{ij}(n)=\int_0^1z^{n-1}P_{ij}(z)\,\dd z$（第~8~章）。与其他文献的约定换算：Moch--Vermaseren--Vogt 的 $\gamma^{(0)}_{\rm MVV}(n)=-2\gamma(n)$（展开参数 $\as/(4\pi)$）；Gross--Wilczek 的 $\mu\,\dd O/\dd\mu=-\gamma_OO$ 对应 $\gamma_O=-\frac{\as}\pi\gamma_{qq}(n)=\frac{\as}{2\pi}C_F\big[1-\frac2{n(n+1)}+4\sum_{j=2}^n\frac1j\big]$。
\par
非微扰地位：强子矩阵元 $A^{(n)}$：\cref{asm:intro-wightman}、\cref{asm:intro-composite}（MA0）。
\uses{prop:dis-ope-euclid,prop:dis-ope-gi,prop:dis-ope-forward,thm:dis-ope,thm:prelim-composite-renorm,prop:appJ-local-operators,prop:qcd-rge,lem:appM-uniqueness,asm:intro-wightman,asm:intro-composite}
\end{proposition}

\begin{proof}
完整证明见附录~M 的\cref{sec:appM-brst}，它直接从\cref{prop:dis-ope-euclid}的 (A3) 推出，不需要系数的唯一性与条件 (NE)。(1) 令 $\hat{\mathcal D}=\mu\partial_\mu+\beta(\as,\varepsilon)\partial_{\as}+\beta_\xi\partial_\xi$，它是在固定 $(g_0,\xi_0)$ 下对 $\mu$ 的全导数，故 $\hat{\mathcal D}C^{\rm bare}=0$；由 $C^{\overline{\rm MS}}=C^{\rm bare}Z^{-1}$ 得 $\hat{\mathcal D}C^{\overline{\rm MS}}=C^{\overline{\rm MS}}\gamma$，$\gamma=-(\hat{\mathcal D}Z)Z^{-1}$，这在完整的 $D$ 维基（含 evanescent 算符）中成立。$\gamma$ 有限：在 $\overline{\rm MS}$ 中 $\gamma=2\as\partial_{\as}Z_1$ 加上极点，而最高阶极点若不为零将与\cref{lem:appM-uniqueness}(c) 矛盾。$\beta_\xi$ 一项是必要的：$\beta_\xi=\mu\,\dd\xi/\dd\mu=\frac{\as}{2\pi}\xi\big[(\frac{13}6-\frac\xi2)C_A-\frac43T_Fn_f\big]+\order{\as^2}$，在 $\xi=1$、SU(3) 时为 $\frac{\as}{2\pi}(5-\frac23n_f)\ne0$（\texttt{s4\_wilson\_rg.py}）。(2) 在 $\xi=1$ 把方程限制到扭度 2 规范不变算符：规范不变与外来算符之外的系数为零（\cref{prop:dis-ope-gi}），外来算符不混入规范不变算符且 $\partial_\xi\gamma_{GG}=0$（\cref{prop:appJ-local-operators}、\cref{thm:prelim-composite-renorm}），扭度 2 系数与 $\xi$ 无关（\cref{prop:dis-ope-forward}：$T^{(j)}_2$ 是在壳物理振幅，$A^{(n)}=Z_{GG}\cdot{\rm tree}$ 与 $\xi$ 无关），于是 $(\mu\partial_\mu+\beta\partial_{\as})C^{GI}_i=\sum_{j\in GI}C^{GI}_j(\gamma_{GG})_{ji}$。(3) 换成 $\ln\mu^2$：$\mu\,\dd/\dd\mu=2\,\dd/\dd\ln\mu^2$，令 $\gamma=-2\Gamma$ 即得\cref{eq:dis-rg}的第二式。裸算符乘积的矩阵元与 $\mu$ 无关，$\sum_iC_iA_i$ 不依赖于 $\mu$，故 $A$ 满足第一式；对强子矩阵元这是\cref{asm:intro-composite}(C1) 的内容，而重整化群方程本身来自\cref{prop:qcd-rge}。(4) 最低阶的 $\gamma_{ij}(n)$ 是单圈的算符反常维度；$\gamma_{qq}$、$\gamma_{qg}$ 在\cref{sec:dis-wilson}的单圈计算中作为 $1/\varepsilon$ 极点的留数重新出现。换算表直接由定义比较得到。
\end{proof}

偶数 $n$ 与奇数 $n$ 的算符分属 C 偶（$q+\bar q$）与 C 奇（$q-\bar q$）的组合；由偶数 $n$ 与由奇数 $n$ 出发的解析延拓在次领头阶以上给出不同的非单态核 $P^\pm_{\rm NS}$（\cref{rem:dis-uniqueness}，第~8~章\cref{rem:pdf-ns-plusminus}）。
